\documentclass[10pt,a4paper]{article}

% ----------------- Paquetes -------------------%
\usepackage[T1]{fontenc}
\usepackage[utf8]{inputenc}
\usepackage[a4paper,margin=2.5cm]{geometry}
\usepackage{mathptmx}             
\usepackage{changepage} 
\usepackage{setspace}
\usepackage[spanish]{babel}
\usepackage[labelfont=bf,labelsep=space]{caption}
\usepackage{amssymb}
\usepackage{amsmath}
\usepackage{ebgaramond}
\usepackage{placeins}
\usepackage{etoolbox}
\usepackage{setspace}
\usepackage{parskip} 
\usepackage{tcolorbox}
\usepackage{needspace}
\usepackage{xparse}
\usepackage{ocgx2}
\usepackage{hyperref}
\usepackage{hypcap}


%%%%%%%%%%%%%%%%%%%%%%%%%%%%%%%%%%%%%%%%%%%%%%%%%%%%%%%%%%%%%%%%
% --- Fija primero tus valores globales "buenos" ---
\setstretch{1.2}                 % Interlineado global
\setlength{\parskip}{0.7\baselineskip}
\setlength{\parindent}{0pt}
\newlength{\GlobalParskip}
\setlength{\GlobalParskip}{\parskip}

\newcounter{eqnum}[subsection]
\renewcommand{\theeqnum}{\arabic{eqnum}}
\newcounter{alphaitem}
\newcounter{numitem}

%% ------------------- Recuadros----------------------%%
\newcommand{\puntosclave}[1]{%
  \begin{tcolorbox}[
    colframe=black,
    title={Puntos clave},
    before upper={\setlength{\parindent}{0.5cm}\setlength{\parskip}{0.5\baselineskip}}
  ]
  #1
  \end{tcolorbox}
}

\newcommand{\modificaciones}[1]{
  \begin{tcolorbox}[
    colback=white,
    colframe=red,
    title={Modificaciones a realizar},
    before upper={\setlength{\parindent}{0.5cm}\setlength{\parskip}{0.5\baselineskip}}
  ]
  #1
  \end{tcolorbox}
}

%% ------------------- Listas----------------------%%
    % --- Lista alfabética ---
    \newenvironment{la}
      {\begin{list}{\alph{alphaitem})}{%
          \usecounter{alphaitem}%
          \setlength{\leftmargin}{1cm}%
          \setlength{\parskip}{3pt}% 
          \setlength{\itemsep}{0pt}% ← espacio entre ítems
          \setlength{\parsep}{0pt}%  ← párrafos dentro de un ítem
      }}
      {\end{list}}
      
    % --- Lista numérica ---
    \newenvironment{lnum}
      {\begin{list}{\arabic{numitem}.}{%
          \usecounter{numitem}%
          \setlength{\leftmargin}{1cm}%
          \setlength{\parskip}{3pt}% 
          \setlength{\itemsep}{0pt}% ← espacio entre ítems
          \setlength{\parsep}{0pt}%  ← párrafos dentro de un ítem
      }}
      {\end{list}}
      
% ------------------- Imágenes----------------------%
    \graphicspath{{Img/}}% Rutas por defecto 
    % --- Imagen adaptativa ---
    % Registros de dimensión definidos una sola vez
    \newdimen\imgcap
    \newdimen\imgmin
    \newdimen\imgmax
    \newdimen\imgremain
    \newdimen\imgh
    \newdimen\imgneed
    
    \newcommand{\imagenfit}[3]{%
      \addvspace{10pt}% separación antes
      \par\begingroup
        % Parámetros de composición (asignaciones, no \newdimen)
        \imgcap=1\baselineskip            % alto estimado de título+fuente
        \imgmin=0.15\textheight
        \imgmax=0.15\textheight
        \imgremain=\pagegoal
        \ifdim\pagegoal=\maxdimen \imgremain=0pt\fi
        \advance\imgremain by -\pagetotal
        \advance\imgremain by -\imgcap
        % Decisión de altura destino
        \ifdim\imgremain<\imgmin
          \newpage
          \imgh=0.20\textheight
        \else
          \imgh=\imgremain
          \ifdim\imgh>\imgmax \imgh=\imgmax \fi
        \fi
        % Reservar espacio para evitar cortes del bloque completo
        \imgneed=\imgh \advance\imgneed by \imgcap
        \Needspace{\imgneed}
        % Cuerpo
        \noindent\begin{minipage}{\textwidth}
          \centering
          \captionof{figure}{\textemdash\ #2}\par\vspace{0.3em}% título arriba
          \includegraphics[width=\linewidth,height=\imgh,keepaspectratio]{\detokenize{#1}}\par
          \vspace{0.3em}\small\textit{Fuente: }#3% fuente abajo
        \end{minipage}%
      \endgroup
      \par\addvspace{10pt}%
    }

%------------ Bloque de RECUADRO ESPECIAL -------------%
    \newenvironment{specialblock}[1]{%
      \begin{quote} % crea sangría a izquierda y derecha
      \small % texto más pequeño
      \noindent\textbf{#1}\\[-0.3em] % título en negrita
      \rule{\linewidth}{0.4pt}\\[0.6em]}{
      \end{quote}}
      

%-------------------------- Author Font -----------------------%
    \newcommand{\authorfont}[1]{{\fontfamily{EBGaramond-TLF}\selectfont #1}}

%-------------------------- Anexos -----------------------%
    \makeatletter
    \newcounter{anexo}
    \renewcommand{\theanexo}{\Roman{anexo}}
    \newcommand{\anexo}[1]{%
      \clearpage
      \phantomsection
      \refstepcounter{anexo}%
      \noindent\textbf{Anexo \theanexo.- }#1\par\medskip
      \addcontentsline{stc}{section}{Anexo \theanexo.- #1}%
    }
    \makeatother

    %-------------------------- Lista de figuras (personal) -----------------------%
\makeatletter
\newcommand{\MyListOfFigures}{%
  \clearpage
  \phantomsection
  {\centering\bfseries\Large Índice de figuras\par}%
  \medskip
  % Entrada en nuestro índice de contenido (stc)
  \addcontentsline{stc}{section}{Índice de figuras}%
  % Recuperación del listado (sin cabecera propia)
  \begingroup
    \parindent=0pt\parskip=2pt
    \@starttoc{lof}%
  \endgroup
}
\makeatother

%-------------------------- Bibliografía -----------------------%
\makeatletter
% Contador para las referencias
\newcounter{myref}
% Cabecera de la bibliografía, entra en el índice 'stc'
\newcommand{\MyBibliography}{%
  \clearpage
  \phantomsection
  {\centering\bfseries\Large Bibliografía y referencias\par}%
  \medskip
  \addcontentsline{stc}{section}{Bibliografía y referencias}%
  \setcounter{myref}{0}%
}
\newcommand{\MyBibItem}[4]{%
  \refstepcounter{myref}%
  \noindent[\themyref]\ #1.\ \emph{#2}. #3. #4\par
}
\makeatother

%--------- Establecer cuatro niveles de índice --------%
    \setcounter{tocdepth}{4}   
    \setcounter{secnumdepth}{4}% numera hasta \paragraph

%%%%%%%%%%%%%%%%%%%%%%%%%%%%%%%%

\makeatletter

    %--------- Interlineado Global --------%
    \edef\GlobalBaseLineStretch{\baselinestretch}
    
    %--------- Espaciado de seccinones --------%
    \renewcommand\section{\@startsection{section}
      {1}{\z@}
      {2.5ex \@plus 1ex \@minus .2ex} % espacio antes
      {1.5ex \@plus .2ex}             % espacio después
      {\normalfont\Large\bfseries}    % formato del título
    }
    \renewcommand\subsection{\@startsection{subsection}{2}{\z@}
      {3ex \@plus 0.8ex \@minus .2ex}
      {1.5ex \@plus .2ex}
      {\normalfont\large\bfseries}
    }
    \renewcommand\subsubsection{\@startsection{subsubsection}{3}{\z@}
      {3ex \@plus 0.5ex \@minus .2ex}
      {1ex \@plus .2ex}
      {\normalfont\normalsize\bfseries}
    }
    \renewcommand\paragraph{\@startsection{paragraph}{4}{\z@}%
    {-2ex \@plus -0.5ex \@minus -.2ex}% espacio antes
    {0.8ex \@plus .2ex}% espacio después
    {\normalfont\normalsize\itshape}}% ← cursiva en lugar de negrita

    %--------- Ecuaciones --------%   
    % Variables globales para almacenar títulos y notas
    \gdef\leq@current@section{}%
    \gdef\leq@current@subsection{}%
    \gdef\leq@last@printed@section{}%
    \gdef\leq@last@printed@subsection{}%
    
    % Comando para añadir nota (invisible en el cuerpo)
    \newcommand{\eqnote}[1]{%
      \addcontentsline{leq}{leqnote}{#1}%
    }
    
    % Comando para ecuaciones
    \newcommand{\eqblock}[2]{%
      % Verificar si necesitamos imprimir el título de sección
      \ifx\leq@current@section\leq@last@printed@section\else
        \addcontentsline{leq}{leqsection}{\leq@current@section}%
        \global\let\leq@last@printed@section\leq@current@section
        \gdef\leq@last@printed@subsection{}% Reset subsección al cambiar sección
      \fi
      % Verificar si necesitamos imprimir el título de subsección
      \ifx\leq@current@subsection\leq@last@printed@subsection\else
        \ifx\leq@current@subsection\@empty\else
          \addcontentsline{leq}{leqsubsection}{\leq@current@subsection}%
          \global\let\leq@last@printed@subsection\leq@current@subsection
        \fi
      \fi
      % Ecuación
      \refstepcounter{eqnum}%
      \begin{equation*}
        #1\tag{E.\theeqnum}
      \end{equation*}%
      \addcontentsline{leq}{leqt}{[E.\theeqnum] -- #2}%
      \addcontentsline{leq}{leqm}{#1}%
    }
    
    %%% ----- Índice de ecuaciones ----------%%%
    % Tamaño para []
    \newcommand{\leqIndexFont}{\normalsize}
    \newcommand{\leqTitleFont}{\tiny}
    \newcommand{\leqNoteFont}{\tiny}
    % Cabecera de sección 
    \newcommand*\l@leqsection[2]{%
      \addvspace{25pt}% Mayor separación antes de nueva sección
      \noindent\leqIndexFont
      \makebox[\linewidth]{\rule{.38\linewidth}{0.4pt}\ \ \textbf{  \MakeUppercase{#1}}\ \ \rule{.38\linewidth}{0.4pt}}%
      \par\vspace{-10pt}
    }
    % Cabecera de subsección 
    \newcommand*\l@leqsubsection[2]{%
      \par\addvspace{15pt}%
      \noindent\leqIndexFont #1
      \par\vspace{-7pt}
    }
    % Título de ecuación 
    \newcommand*\l@leqt[2]{%
    \par\addvspace{10pt}%
      \par\noindent\leqTitleFont #1\par
    }
    % Miniatura de ecuación
    \newcommand*\l@leqm[2]{%
      \vspace{0.3\baselineskip}%
      \par\scriptsize\noindent\hspace*{1.5em}\(#1\)\par%
    }
    % Nota de ecuación 
    \newcommand*\l@leqnote[2]{%
      \vspace{0.2\baselineskip}%
      \par\noindent\leqNoteFont\hspace*{1.5em}\textbf{Nota.--} #1\par\vspace{0.5\baselineskip}%
    }
    % Generación del índice
    \newcommand{\mylistofequations}{%
      \clearpage
      \phantomsection
      % Título visual de la lista (ajústalo a tu gusto):
      {\centering\bfseries\Large Índice de ecuaciones\par}
      % Entrada en nuestro índice 'stc':
      \addcontentsline{stc}{section}{Índice de ecuaciones}%
      % Llamada a tu lista real (usa el nombre exacto de tu macro):
      \@starttoc{leq}%
    }
    
    %--------- Captura de títulos de sección --------%
    \let\leq@original@section\section
    \let\leq@original@subsection\subsection
    % Flag para saber si estamos en una sección asterisco
    \newif\ifLeqStarSection
    \LeqStarSectionfalse
    
    % Redefinir \section CON CONTROL DE ASTERISCO
    \renewcommand{\section}{%
      \@ifstar{\leq@section@star}{\@dblarg\leq@section@nostar}%
    }
    \def\leq@section@star#1{%
      \global\LeqStarSectiontrue
      \gdef\leq@current@section{#1}%
      \gdef\leq@current@subsection{}%
      \leq@original@section*{#1}%
      \global\LeqStarSectionfalse
    }
    \def\leq@section@nostar[#1]#2{%
      \global\LeqStarSectionfalse
      \gdef\leq@current@section{#2}%
      \gdef\leq@current@subsection{}%
      \leq@original@section[#1]{#2}%
    }
    % Redefinir \subsection
    \renewcommand{\subsection}{%
      \@ifstar{\leq@subsection@star}{\@dblarg\leq@subsection@nostar}%
    }
    \def\leq@subsection@star#1{%
      \global\LeqStarSectiontrue
      \gdef\leq@current@subsection{#1}%
      \leq@original@subsection*{#1}%
      \global\LeqStarSectionfalse
    }
    \def\leq@subsection@nostar[#1]#2{%
      \global\LeqStarSectionfalse
      \gdef\leq@current@subsection{#2}%
      \leq@original@subsection[#1]{#2}%
    }

    % ----- Restauración de valores Globales de estilo ----------%
    % Flags
    \newif\ifInFootnote
    \InFootnotefalse
    \pretocmd{\@footnotetext}{\InFootnotetrue}{}{}
    \apptocmd{\@footnotetext}{\InFootnotefalse}{}{}
    % Profundidad de listas (soporta anidadas)
    \newcount\MyListDepth \MyListDepth=0
    \AtBeginEnvironment{la}{\advance\MyListDepth by 1}
    \AfterEndEnvironment{la}{\advance\MyListDepth by -1}
    \AtBeginEnvironment{lnum}{\advance\MyListDepth by 1}
    \AfterEndEnvironment{lnum}{\advance\MyListDepth by -1}
    \AtBeginEnvironment{itemize}{\advance\MyListDepth by 1}
    \AfterEndEnvironment{itemize}{\advance\MyListDepth by -1}
    \AtBeginEnvironment{enumerate}{\advance\MyListDepth by 1}
    \AfterEndEnvironment{enumerate}{\advance\MyListDepth by -1}
    % Restauración diferida: hazlo al INICIO del próximo párrafo real
    \newif\if@pendingRestore
    \@pendingRestorefalse
    \newcommand{\RequestRestore}{\global\@pendingRestoretrue}
    \everypar{%
      \if@pendingRestore
        \global\@pendingRestorefalse
        \ifInFootnote\else
          \ifnum\MyListDepth=0
            \setlength{\parskip}{\GlobalParskip}%
            \setstretch{\GlobalBaseLineStretch}%
          \fi
        \fi
      \fi
    }
    % Al final de bloques:
    \AfterEndEnvironment{equation}{\RequestRestore}
    \AfterEndEnvironment{equation*}{\RequestRestore}
    \AfterEndEnvironment{la}{\RequestRestore}
    \AfterEndEnvironment{lnum}{\RequestRestore}
    % Al final de macros:
    \apptocmd{\eqblock}{\RequestRestore}{}{}
    \apptocmd{\imagenfit}{\RequestRestore}{}{}
    
\makeatother

% ===================== ToC propio con 4 niveles (único bloque) =====================
\makeatletter

% --- Escritores: añadimos una línea al .stc en cada cabecera numerada ---
% Guardamos las versiones actuales para llamar al comportamiento normal
\let\MyTOC@orig@section\section
\let\MyTOC@orig@subsection\subsection
\let\MyTOC@orig@subsubsection\subsubsection
\let\MyTOC@orig@paragraph\paragraph

% SECTION
\renewcommand{\section}{%
  \@ifstar{\MyTOC@section@star}{\@dblarg\MyTOC@section@nostar}%
}
\def\MyTOC@section@star#1{%
  \MyTOC@orig@section*{#1}% sin entrada al .stc
}
\def\MyTOC@section@nostar[#1]#2{%
  \MyTOC@orig@section[{#1}]{#2}% comportamiento normal (escribe al .toc)
  % Escribimos también nuestra línea al .stc (texto con \numberline)
  \addcontentsline{stc}{section}{\protect\numberline{\thesection}#2}%
}

% SUBSECTION
\renewcommand{\subsection}{%
  \@ifstar{\MyTOC@subsection@star}{\@dblarg\MyTOC@subsection@nostar}%
}
\def\MyTOC@subsection@star#1{%
  \MyTOC@orig@subsection*{#1}%
}
\def\MyTOC@subsection@nostar[#1]#2{%
  \MyTOC@orig@subsection[{#1}]{#2}%
  \addcontentsline{stc}{subsection}{\protect\numberline{\thesubsection}#2}%
}

% SUBSUBSECTION
\renewcommand{\subsubsection}{%
  \@ifstar{\MyTOC@subsubsection@star}{\@dblarg\MyTOC@subsubsection@nostar}%
}
\def\MyTOC@subsubsection@star#1{%
  \MyTOC@orig@subsubsection*{#1}%
}
\def\MyTOC@subsubsection@nostar[#1]#2{%
  \MyTOC@orig@subsubsection[{#1}]{#2}%
  \addcontentsline{stc}{subsubsection}{\protect\numberline{\thesubsubsection}#2}%
}

% PARAGRAPH
\renewcommand{\paragraph}{%
  \@ifstar{\MyTOC@paragraph@star}{\@dblarg\MyTOC@paragraph@nostar}%
}
\def\MyTOC@paragraph@star#1{%
  \MyTOC@orig@paragraph*{#1}%
}
\def\MyTOC@paragraph@nostar[#1]#2{%
  \MyTOC@orig@paragraph[{#1}]{#2}%
  % Si no numeras los \paragraph, cambia \theparagraph por texto plano sin \numberline
  \addcontentsline{stc}{paragraph}{\protect\numberline{\theparagraph}#2}%
}

% --- Impresor del índice propio (lee el .stc) ---
\newcommand{\MyTOC}{%
  \par\bigskip
  {\centering\bfseries\Large Índice de contenido\par}%
  \medskip
  \begingroup
    \parindent=0pt\parskip=2pt
    \@starttoc{stc}%
  \endgroup
}

\makeatother
% ===================== fin ToC propio =====================


%%%%%%%%%%%%%%


% --- Datos del tema ---
\title{El presupuesto como elemento compensador de la actividad económica\\
\large Componentes discrecionales y automáticos del presupuesto. Saldo cíclico, saldo estructural y esfuerzo fiscal}
\author{Marcelo Ortega}
\date{Septiembre 2026}

\usepackage{soul}\sethlcolor{yellow}% resaltado amarillo: \hl{texto} (Panel TCEE, ⌃H)
\begin{document}
\maketitle
\newpage

% --- Página de resumen y modificaciones ---
\puntosclave{ %% Escribir puntos clave Apuntes CECO
     
    }
\vspace{1cm}

\modificaciones{
    
    }

\newpage
\MyTOC
\newpage

\mylistofequations
\clearpage

\MyListOfFigures
\clearpage


\pagenumbering{arabic}
\setcounter{page}{1}


\section{Introducción}



\subsection{Contextualización}


\subsubsection{Relevancia}




\subsubsection{Problemática}



\subsection{Estructura de la exposición}





\subsubsection{\textbf{AQUI!}}

\textcolor{red}{\textbf{ORIGEN:} ROSEN y GARBER (Capítulo 12)}

Incidencia del gasto

Pasamos ahora de una discusión sobre si el gobierno debería redistribuir la renta a los problemas analíticos de evaluar los efectos de los programas redistributivos reales del gobierno. El impacto de la política de gasto sobre la distribución de la renta real se denomina incidencia del gasto. El gobierno influye en la distribución de la renta tanto a través de sus políticas tributarias como de sus políticas de gasto. (Posponemos una discusión del lado tributario hasta el capítulo 14.) La incidencia del gasto es difícil de determinar por varias razones, que se exponen a continuación.

Efectos sobre los precios relativos

Supongamos que el gobierno decide subvencionar el consumo de vivienda para personas de bajos ingresos. ¿Cómo afecta esto a la distribución de la renta? Una primera conjetura sería que las personas que reciben la subvención ganan y aquellas que pagan los impuestos pierden. Si quienes pagan los impuestos tienen rentas más altas que los receptores de la subvención, la distribución de la renta se vuelve más igualitaria.

Desafortunadamente, esta sencilla historia puede ser engañosa. Si la subvención induce a las personas pobres a demandar más vivienda, entonces el coste de la vivienda anterior a la subvención puede aumentar. Por lo tanto, los receptores de la subvención no se benefician en toda la cuantía de la subvención; los propietarios se llevan parte de la ganancia. Sin embargo, únicamente sobre bases teóricas no puede determinarse cuánto, si es que algo, aumentan los precios de la vivienda. Como se muestra en el capítulo 14, esto depende de las formas de las curvas de oferta y demanda de vivienda.

Un programa de subvenciones a la vivienda también afecta a las rentas de las personas que suministran los factores utilizados en su construcción. Así, aumentan los salarios de los trabajadores de los oficios de la construcción, al igual que los precios de los materiales de construcción. Si los propietarios de estos factores pertenecen a las clases media y alta, esto tenderá a hacer que la distribución sea menos igualitaria.

De manera más general, cualquier programa gubernamental pone en marcha una cadena de cambios de precios que afecta a las rentas de las personas tanto en sus papeles como consumidores de bienes como en sus papeles como proveedores de factores. Un programa de gasto que eleva el precio relativo de un bien que usted consume hace que usted esté peor, permaneciendo iguales las demás cosas. Del mismo modo, un programa que eleva el precio relativo de un factor que usted suministra hace que usted esté mejor. El problema es que resulta muy difícil rastrear todos los cambios de precios generados por una política determinada. Como cuestión práctica, los economistas suelen suponer que una política determinada beneficia únicamente a los receptores y que los efectos de otros cambios de precios sobre la distribución de la renta son menores. En muchos casos, probablemente sea un buen supuesto.

Bienes públicos

Una parte sustancial del gasto gubernamental se destina a bienes públicos —bienes que pueden ser consumidos simultáneamente por más de una persona—. Como se señaló en el capítulo 4, el mercado no obliga a las personas a revelar cuánto valoran los bienes públicos. Pero si no sabemos cuánto valora cada familia un bien público, ¿cómo podemos determinar su impacto sobre la distribución de la renta? El gobierno gastó más de 650.000 millones de dólares en defensa en 2008. ¿En cuánto, en términos de dólares, aumentó esto la renta real de cada familia? ¿Se benefició cada una en la misma cantidad? Si no, ¿se beneficiaron los pobres menos que los ricos, o viceversa?

Es imposible responder definitivamente a preguntas como estas. Desafortunadamente, respuestas alternativas basadas en supuestos igualmente plausibles tienen implicaciones muy diferentes. Chamberlain y Prante [2007] examinaron las implicaciones distributivas de los gastos en bienes públicos como la defensa utilizando dos supuestos diferentes: (a) la participación de un hogar en el beneficio es proporcional a su riqueza, y (b) cada hogar recibe una participación igual del beneficio. Bajo el supuesto (a), el quinto superior de la población recibe el 27,0 por ciento del gasto total del gobierno en bienes públicos, mientras que bajo el supuesto (b) recibe solo el 17,1 por ciento de los gastos en bienes públicos. Los resultados son claramente muy sensibles a los supuestos.

Valoración de las transferencias en especie

Durante las últimas varias décadas, el Departamento de Agricultura ha regalado más de 3.000 millones de libras de excedentes de queso, mantequilla y leche en polvo a estadounidenses pobres. El programa de excedentes de alimentos es solo un ejemplo de una política de transferencias en especie. A menudo pensamos en las transferencias en especie como dirigidas a individuos de renta más baja: vienen a la mente los cupones de alimentos, Medicaid y la vivienda pública. Sin embargo, las personas de renta media y alta también se benefician de las transferencias en especie. Un ejemplo destacado es la educación.

A diferencia de los bienes públicos puros, las transferencias en especie no son consumidas por todo el mundo. No obstante, estimar su valor para los beneficiarios es difícil. Un supuesto conveniente es que un dólar gastado por el gobierno en una transferencia en especie equivale a un aumento de un dólar en la renta del receptor. Desafortunadamente, no hay ninguna razón para creer que las transferencias en especie sean valoradas por los beneficiarios dólar por dólar.

Para ver por qué, considérese a Jones, una receptora típica de asistencia social que divide su renta mensual de 300 dólares entre queso y «todos los demás bienes». El precio de mercado del queso es de 2 dólares por libra, y las unidades de «todos los demás bienes» se miden de modo que el precio por unidad sea de 1 dólar. En la Figura 12.3, el consumo de queso de Jones se mide en el eje horizontal, y su consumo de todos los demás bienes en el vertical. La restricción presupuestaria de Jones es la línea AB.⁶ Suponiendo que Jones maximiza su utilidad, consume la cesta E₁, que consta de 260 unidades de todos los demás bienes y 20 libras de queso.

Supongamos ahora que el gobierno proporciona a Jones 60 libras de queso al mes, que tiene prohibido revender en el mercado. ¿Cómo cambia su situación la introducción del programa de queso? Para cualquier nivel de consumo de todos los demás bienes, Jones puede ahora consumir 60 libras más de queso que anteriormente. Geométricamente, su nueva restricción presupuestaria se obtiene desplazando 60 unidades hacia la derecha cada punto de AB, obteniendo AFD. La curva de indiferencia más alta que puede alcanzar sujeta a la restricción AFD es la curva U de la Figura 12.3. Esta toca la restricción en su «esquina» —en el punto F, donde el consumo de queso de Jones es 60 y su consumo de todos los demás bienes es 300—.

En comparación con su cesta de consumo original, el consumo de Jones tanto de queso como de todos los demás bienes ha aumentado. Debido a que el gobierno le proporciona queso gratis, Jones puede utilizar el dinero que habría gastado en queso para comprar más de todos los demás bienes.

Supongamos ahora que, en lugar de dar a Jones 60 libras de queso, el gobierno le da efectivo igual a su valor de mercado, 120 dólares (60 libras × 2 dólares por libra). Un aumento de la renta de 120 dólares conduce a una línea presupuestaria que está exactamente 120 unidades por encima de AB en cada punto, representada en la Figura 12.3 como la línea HD. Obsérvese que la transferencia en efectivo permite a Jones consumir a lo largo del segmento HF. Esta oportunidad no estaba disponible bajo el programa de queso porque a Jones no se le permitía intercambiar el queso del gobierno por ningún otro bien.

Frente a la línea presupuestaria HD, Jones maximiza la utilidad en el punto E₃, donde consume 340 de todos los demás bienes y 40 libras de queso. Comparando los puntos E₃ y F podemos concluir que (1) bajo el programa de transferencias en efectivo, Jones consume menos queso y más de todos los demás bienes que bajo el programa de entrega gratuita de queso; y (2) 120 dólares de queso no hacen que Jones esté tan bien como 120 dólares de renta. Debido a que E₃ se encuentra en una curva de indiferencia más alta que el punto F, la transferencia en efectivo hace que esté mejor.

Intuitivamente, el problema con el programa de queso es que obliga a Jones a consumir las 60 libras completas de queso. Ella preferiría vender parte del queso y gastar los ingresos en otros bienes.

¿Es una transferencia en especie siempre peor que su equivalente en efectivo? No necesariamente. La Figura 12.4 representa la situación de Smith, cuya renta es idéntica a la de Jones y que, por tanto, se enfrenta exactamente a las mismas restricciones presupuestarias (AB antes del programa de queso y AFD después). Sin embargo, Smith tiene gustos diferentes y, por tanto, un conjunto diferente de curvas de indiferencia. Antes de la subvención, maximiza la utilidad en el punto E₄, consumiendo 136 unidades de todos los demás bienes y 82 libras de queso. Después de la subvención, consume 168 unidades de todos los demás bienes y 126 libras de queso. Smith no estaría mejor con una transferencia en efectivo porque su punto más preferido a lo largo de HD está disponible de todos modos bajo la subvención de queso. Debido a que Smith está contento de consumir más de 60 libras de queso, la restricción de que consuma al menos 60 libras no le causa ningún perjuicio.

Así, no podemos saber con certeza si una transferencia en especie se valora menos que una transferencia directa de renta. En última instancia, la respuesta tiene que encontrarse mediante análisis empírico. Por ejemplo, un estudio estima que un dólar recibido en cupones de alimentos (vales que solo pueden utilizarse para comprar alimentos) vale solo alrededor de 80 centavos recibidos en efectivo [Whitmore, 2002].

Otro problema de los programas de transferencias en especie es que a menudo implican costes administrativos sustanciales. En el programa de queso que acabamos de analizar, se incurre en costes por el almacenamiento, transporte y distribución del queso. (Los costes son tan elevados que algunas comunidades optan por no participar.) Del mismo modo, los costes administrativos del programa de cupones de alimentos podrían reducirse si los beneficiarios simplemente recibieran cheques en lugar de cupones canjeables por alimentos.

Razones para las transferencias en especie

Como mostramos en el capítulo 13, las transferencias en especie que implican alimentos, vivienda y atención médica desempeñan un papel importante en la política estadounidense de mantenimiento de la renta. Si las transferencias en especie son menos beneficiosas que el efectivo desde el punto de vista de los receptores e implican más costes administrativos, ¿cómo podemos explicar su presencia? Hay una serie de posibles explicaciones. Varias están relacionadas con nuestra discusión anterior de las cuestiones normativas.

En particular, el igualitarismo de bienes puede desempeñar un papel importante en la política distributiva. Por ejemplo, el Congreso de Estados Unidos estableció una vez explícitamente como objetivo nacional «una vivienda digna y un entorno de vida adecuado para cada familia estadounidense». Obsérvese la distinción entre este objetivo y «renta suficiente para que cada familia estadounidense pueda vivir en una vivienda digna, si así lo elige».

Además, las transferencias en especie también pueden ayudar a frenar el fraude en la asistencia social. La discusión hasta ahora ha supuesto que no existen problemas para identificar quién reúne los requisitos para recibir una transferencia y quién no. En realidad, este no es el caso, y las personas que no reúnen los requisitos a veces pueden obtener prestaciones. Las transferencias en especie pueden disuadir a las personas no elegibles de solicitarlas porque algunas personas de clase media pueden estar bastante dispuestas a mentir para recibir efectivo, pero menos dispuestas a mentir para obtener un bien que realmente no quieren. Esto es especialmente cierto si el bien es difícil de revender, como un apartamento en un proyecto de vivienda pública. Del mismo modo, crear molestias para los receptores de asistencia social (esperar en fila, rellenar muchos formularios) puede disuadir de solicitarla a aquellos que no están «verdaderamente necesitados».

Así, existe una disyuntiva. Por una parte, una persona pobre preferiría 500 dólares en efectivo a 500 dólares en vivienda pública. Pero si el programa en especie conduce a menos fraude, pueden canalizarse más recursos hacia las personas que realmente los necesitan. Sin embargo, muchos sostendrían que el gobierno ha creado muchas más barreras administrativas que el número óptimo para los receptores de asistencia social. Por ejemplo, en 2003 la administración Bush propuso que, para recibir almuerzos escolares gratuitos, los estudiantes tuvieran que proporcionar pruebas, como recibos de sueldo, de que las rentas de sus padres eran suficientemente bajas. Algunos observadores consideraron esto una carga injusta para los niños.

Finalmente, las transferencias en especie son políticamente atractivas porque ayudan no solo al beneficiario, sino también a los productores del bien favorecido. Un programa de transferencias que aumenta la demanda de vivienda beneficia a la industria de la construcción, que, por tanto, está dispuesta a prestar su apoyo a una coalición política a favor del programa. Del mismo modo, los intereses agrícolas siempre han sido ávidos partidarios de los cupones de alimentos. Cuando el estado de Oregón pidió permiso para convertir los cupones de alimentos en efectivo para los receptores de asistencia social hace varios años, la idea fue bloqueada por miembros del Congreso de estados agrícolas. Del mismo modo, los empleados públicos que administran los diversos programas de transferencias en especie ponen su apoyo político detrás de ellos. Por ejemplo, los burócratas del Departamento de Vivienda y Desarrollo Urbano se han opuesto tradicionalmente de manera enérgica a las propuestas de que la vivienda subvencionada sea eliminada gradualmente y sustituida por subvenciones en efectivo.

Estas explicaciones de las transferencias en especie no son mutuamente excluyentes, y probablemente todas ellas han influido en el diseño de las políticas.

Conclusión

Hemos examinado una amplia gama de opiniones relativas a la conveniencia de políticas gubernamentales explícitas para redistribuir la renta. Las opiniones abarcan toda la gama, desde diseñar una igualdad completa hasta no hacer nada. El alcance del desacuerdo no es sorprendente.

Establecer un objetivo distributivo no es menos que formalizar las propias opiniones sobre cómo debería ser una buena sociedad, y esto está destinado a ser controvertido. Las teorías sobre la distribución óptima de la renta son normativas en lugar de positivas. Como veremos en el capítulo 13, no está claro si alguna teoría normativa coherente es compatible con las prácticas reales de distribución de la renta de Estados Unidos.








\clearpage
\section{Conclusión} 


\subsection{Recapitulación}


\subsection{Extensiones y relación con otras partes del Temario}



\subsection{Opinión}

\subsubsection{Idea final (Salida o cierra)}

\clearpage
\pagenumbering{gobble}
\MyBibliography



\MyBibItem{DAWID y GATTI}{Chapter 2 - Agent-Based Macroeconomics}{2018}{https://doi.org/10.1016/bs.hescom.2018.02.006}


% ANEXOS
\clearpage










\end{document}